\documentclass[10pt,pdf,hyperref={unicode}]{beamer}

\usepackage{fontspec}
\usepackage{booktabs}
\usepackage{polyglossia}
\setmainlanguage{russian}
\setotherlanguage{english}

\newfontfamily\russianfont{FreeSerif}
\newfontfamily\russianfontsf{FreeSans}
\newfontfamily\russianfonttt{FreeMono}

\setsansfont{FreeSans} 
\setmonofont{FreeMono}

\setbeamercolor{author in head/foot}{fg=white, bg=purple!80!black}
\setbeamercolor{title in head/foot}{fg=white, bg=purple!60!black}
\setbeamercolor{date in head/foot}{fg=white, bg=purple!40!black}

\makeatletter
\defbeamertemplate*{footline}{my theme}{
    \leavevmode%
    \hbox{%
    \begin{beamercolorbox}[wd=.3\paperwidth,ht=2.25ex,dp=1ex,center]{author in head/foot}%
        \usebeamerfont{author in head/foot}\insertshortauthor
    \end{beamercolorbox}%
    \begin{beamercolorbox}[wd=.4\paperwidth,ht=2.25ex,dp=1ex,center]{title in head/foot}%
        \usebeamerfont{title in head/foot}\insertshorttitle
    \end{beamercolorbox}%
    \begin{beamercolorbox}[wd=.3\paperwidth,ht=2.25ex,dp=1ex,right]{date in head/foot}%
        \usebeamerfont{date in head/foot}\insertshortdate{}\hspace*{2em}
        \insertframenumber{} / \inserttotalframenumber\hspace*{2ex}
    \end{beamercolorbox}}%
}
\makeatother

\title[Алгоритмы замещения]{Сравнительный анализ алгоритмов замещения страниц}
\author[Макаренко А. А.]{Макаренко Александра Александровна}
\institute{НКАбд-07-25}
\date{2 марта 2026 г.}

\begin{document}

\begin{frame}
    \titlepage
\end{frame}


\begin{frame}{Алгоритм NRU}
    \textbf{NRU (Not Recently Used)} выбирает для замещения страницу, не использовавшуюся в последнее время.
    
    \vspace{0.3cm} 
    
    Используются биты в таблице страниц:
    \begin{itemize}
        \item \textbf{R (Referenced):} 1 — было обращение, 0 — сбрасывается ОС.
        \item \textbf{M (Modified):} 1 — страница изменена, 0 — после записи на диск. 
    \end{itemize}

    \begin{block}{Классы приоритетета на удаление:}
        \begin{enumerate}
            \item \textbf{Класс 0:} R=0, M=0 (не было обращений и изменений)
            \item \textbf{Класс 1:} R=0, M=1 (не было обращений, было изменение)
            \item \textbf{Класс 2:} R=1, M=0 (было обращение, не было изменений)
            \item \textbf{Класс 3:} R=1, M=1 (были обращения и изменения)
        \end{enumerate}
    \end{block}
\end{frame}

\begin{frame}{Алгоритм NRU: Плюсы и минусы}
    \begin{block}{Плюсы}
        \begin{itemize}
            \item \textbf{Эффективность:} быстрая реализация через аппаратные биты.
            \item \textbf{Комплексность:} учитывает и активность, и изменения.
            \item Работает лучше FIFO и \textbf{нет аномалии Биледи}.
        \end{itemize}
    \end{block}

    \begin{alertblock}{Минусы}
        \begin{itemize}
            \item \textbf{Низкая точность:} менее точен, чем LRU.
            \item Требует периодического обнуления бита R.
            \item \textbf{Случайность:} выбор внутри класса случайный $\rightarrow$ возможна неоптимальная замена.
        \end{itemize}
    \end{alertblock}
\end{frame}


\begin{frame}{Алгоритм FIFO}
    \textbf{FIFO (First In, First Out)} — простейший алгоритм: страница, которая зашла в память \textbf{первой}, первой её и \textbf{покидает}.

    \vspace{0.3cm}

    \begin{block}{Плюсы}
        \begin{itemize}
            \item \textbf{Простая реализация:} используется обычная очередь (queue).
            \item \textbf{Экономичность:} практически не требует накладных расходов по памяти.
        \end{itemize}
    \end{block}

    \begin{alertblock}{Минусы}
        \begin{itemize}
            \item \textbf{Низкая эффективность:} часто удаляет страницы, которые активно используются прямо сейчас.
            \item \textbf{Аномалия Биледи:} парадоксальная ситуация, когда увеличение количества кадров памяти приводит к \textit{росту} числа ошибок страниц.
        \end{itemize}
    \end{alertblock}
\end{frame}

\begin{frame}{Алгоритм OPT (Optimal)}
    \textbf{Теоретический идеал:} удаляет страницу, которая потребуется \textbf{позднее всех} в будущем (или не потребуется вовсе).

    \begin{block}{Плюсы}
        \begin{itemize}
            \item Обеспечивает \textbf{минимально возможное} число ошибок страниц (page faults).
            \item Служит «золотым стандартом» (эталоном) для оценки эффективности других алгоритмов.
        \end{itemize}
    \end{block}

    \begin{alertblock}{Минусы}
        \begin{itemize}
            \item \textbf{Нереализуем:} ОС не может предсказать будущее поведение программ и последовательность запросов.
        \end{itemize}
    \end{alertblock}
\end{frame}

\begin{frame}{Алгоритм LRU (Least Recently Used)}
    Основан на принципе локальности: если страница долго не использовалась, вероятно, она не понадобится и в ближайшем будущем.

    \begin{block}{Плюсы}
        \begin{itemize}
            \item Высокая эффективность (по результатам близок к алгоритму OPT).
            \item \textbf{Отсутствие аномалии Биледи} (относится к классу стековых алгоритмов).
        \end{itemize}
    \end{block}

    \begin{alertblock}{Минусы}
        \begin{itemize}
            \item \textbf{Сложность:} требуется постоянное обновление данных при каждом обращении к памяти.
            \item \textbf{Ресурсоемкость:} высокие затраты на поддержку списков или счетчиков времени.
        \end{itemize}
    \end{alertblock}
\end{frame}

\begin{frame}{Алгоритм NFU (Not Frequently Used)}
    Алгоритм подсчитывает частоту обращений: на выход идет та страница, к которой обращались \textbf{реже всего}.

    \begin{block}{Плюсы}
        \begin{itemize}
            \item Проще в реализации, чем LRU (не нужно хранить точный порядок всех обращений).
            \item Меньше нагрузка на систему при каждом доступе.
        \end{itemize}
    \end{block}

    \begin{alertblock}{Минусы}
        \begin{itemize}
            \item \textbf{«Плохая память»:} алгоритм помнит старые обращения. Страница, которая была популярна час назад, будет занимать место, даже если сейчас она не нужна.
            \item Низкая адаптивность к смене фаз работы программы.
        \end{itemize}
    \end{alertblock}
\end{frame}
\begin{frame}{Сравнение алгоритмов замещения страниц}
    \centering
    
    \resizebox{\textwidth}{!}{
        \begin{tabular}{lcccc}
            \toprule
            \textbf{Алгоритм} & \textbf{Реализуемость} & \textbf{Интеллект} & \textbf{Производит.} & \textbf{Аномал.} \\
            \midrule
            OPT  & Нереализуем & Да (будущее) & Лучшее  & Нет \\
            LRU  & Сложная     & Да (прошлое) & Высокая & Нет \\
            NFU  & Средняя     & Частично     & Средняя & Нет \\
            FIFO & Простая     & Нет          & Низкая  & Есть \\
            NRU & Простая     & Частично          & Средняя/выше FIFO  & нет \\
            \bottomrule
        \end{tabular}
    }
\end{frame}


\end{document}
